%% =====================================================================
%%  CO5151 · G4 · SLIDE FINAL — KHUNG (SKELETON)
%%  Tổng 20 phút. Phân bổ theo gợi ý của đề:
%%    2' bài toán/đóng góp  ·  4' kiến trúc + trục nâng cao
%%    6' LIVE DEMO          ·  5' đánh giá + security  ·  3' thay đổi từ D2 + kết luận
%%
%%  QUY TẮC:
%%    - Mỗi frame 1 ý. Tối đa 6 bullet, mỗi bullet <= 12 từ.
%%    - Số trên slide PHẢI khớp báo cáo. Không có số riêng cho slide.
%%    - Frame backup (sau \appendix) KHÔNG tính vào 20 phút.
%%
%%  BIÊN DỊCH:  tectonic main.tex   → main.pdf   (16:9)
%% =====================================================================
\documentclass[aspectratio=169,10pt]{beamer}

\usepackage{fontspec}
\setmainfont{texgyretermes}[Extension=.otf,
  UprightFont=*-regular, BoldFont=*-bold,
  ItalicFont=*-italic, BoldItalicFont=*-bolditalic]
\usepackage[vietnamese]{babel}
\usepackage{booktabs}
\usepackage{graphicx}
\usepackage{xcolor}
\usepackage{tikz}

%% --- Bảng màu (đổi 1 chỗ là đổi cả deck) ---
\definecolor{hcmut}{RGB}{0,80,140}
\definecolor{accent}{RGB}{200,80,0}
\definecolor{ok}{RGB}{0,120,70}
\definecolor{bad}{RGB}{170,30,30}

\usetheme{default}
\usecolortheme{default}
\setbeamercolor{frametitle}{fg=hcmut}
\setbeamercolor{title}{fg=hcmut}
\setbeamercolor{structure}{fg=hcmut}
\setbeamerfont{frametitle}{size=\large,series=\bfseries}
\setbeamertemplate{navigation symbols}{}
\setbeamertemplate{footline}{%
  \hfill\usebeamerfont{footline}\insertframenumber/\inserttotalframenumber\hspace{6pt}\vspace{3pt}}
\setbeamertemplate{blocks}[rounded][shadow=false]

\newcommand{\TODO}[1]{\textcolor{accent}{\textbf{[#1]}}}
\newcommand{\good}[1]{\textcolor{ok}{\textbf{#1}}}
\newcommand{\warn}[1]{\textcolor{bad}{\textbf{#1}}}
\newcommand{\tm}[1]{\hfill{\scriptsize\color{gray}[#1]}}

\title{LegalPilot-VN}
\subtitle{Hệ đa tác tử tư vấn tuân thủ pháp luật Việt Nam}
\author{G4 · Đặng Lâm Tùng · Nguyễn Trung Phong · Vũ Việt Hưng}
\institute{CO5151 · HK261 · 09/12/2026}
\date{}

\begin{document}

%% ---------------------------------------------------------------- 0:00
\begin{frame}
  \titlepage
\end{frame}

%% ------------------------------------------------------- 0:00–0:30
\begin{frame}{Bài toán}
  \begin{itemize}
    \item Chuyên viên tuân thủ ngân hàng: tra \textbf{1.703 văn bản QPPL} bằng tay
    \item Sai 1 dẫn chiếu đã hết hiệu lực $\rightarrow$ rủi ro pháp lý thật
    \item Câu hỏi thực: \emph{“Khoản vay nước ngoài cần điều kiện gì?”}
  \end{itemize}
  \vspace{0.6em}
  \begin{block}{Vì sao phải là agent, không chỉ RAG}
    Phân rã câu hỏi $\rightarrow$ duyệt đồ thị $\rightarrow$ đối soát hiệu lực
    $\rightarrow$ kiểm từng mệnh đề. Bước sau phụ thuộc kết quả bước trước.
  \end{block}
  \tm{0:00–0:30}
\end{frame}

%% ------------------------------------------------------- 0:30–2:00
\begin{frame}{Đóng góp thực tế — chỉ cái đã cài và đã đo}
  \begin{enumerate}
    \item \textbf{Selective Temporal Traversal} — thay 1-hop cố định của [1]
    \item \textbf{Claim Auditor} — kiểm từng mệnh đề nguyên tử với context
    \item \textbf{Citation Gate tất định} — nháp phải trích $\geq1$ căn cứ top-3
    \item \textbf{Bộ đánh giá} 100 câu $\times$ 3 run $\times$ 2 cấu hình, có cost/latency
    \item \textbf{20 ca injection} direct + indirect + API, có log trước/sau
  \end{enumerate}
  \vspace{0.4em}
  \TODO{chốt câu “khác biệt nằm ở cơ chế, không phải ở tên framework”}
  \tm{1:00–2:00}
\end{frame}

%% ------------------------------------------------------- 2:00–3:30
\begin{frame}{Kiến trúc và luồng điều khiển}
  \begin{columns}[T]
    \begin{column}{0.52\textwidth}
      \TODO{\textbf{Hình 1 — sơ đồ THẬT}}
      \vspace{0.3em}
      \begin{enumerate}
        \item cổng guardrail đầu vào
        \item Planner tách 3--5 sub-query
        \item Retriever: Qdrant + Neo4j
        \item Claim Auditor
        \item Drafter + citation gate
      \end{enumerate}
      \TODO{điều kiện dừng + số vòng lặp tối đa}
    \end{column}
    \begin{column}{0.44\textwidth}
      \begin{block}{Điểm khác paper [1]}
        1-hop cố định $\rightarrow$ duyệt \textbf{chọn lọc theo mục tiêu}
      \end{block}
      \TODO{chèn số: max hops, max nodes, top-k}
    \end{column}
  \end{columns}
  \tm{2:00–3:30}
\end{frame}

%% ------------------------------------------------------- 3:30–4:30
\begin{frame}{Tools, phân quyền và bộ nhớ}
  \begin{columns}[T]
    \begin{column}{0.5\textwidth}
      \textbf{Tools} \TODO{bảng thật từ src/tools/}
      \begin{itemize}
        \item đọc: graph, VBPL live, web search
        \item ghi: \warn{chỉ dossier} — cần human token
        \item timeout + fallback mỗi tool
      \end{itemize}
    \end{column}
    \begin{column}{0.46\textwidth}
      \textbf{Memory}
      \begin{itemize}
        \item ngắn hạn: token gate
        \item dài hạn: SQLite session
        \item \textbf{chính sách GHI} \TODO{}
        \item \textbf{chính sách TRUY HỒI} \TODO{}
      \end{itemize}
    \end{column}
  \end{columns}
  \tm{3:30–4:30}
\end{frame}

%% ------------------------------------------------------- 4:30–6:00
\begin{frame}{Ba trục nâng cao ngoài tools/memory}
  \small
  \begin{tabular}{@{}p{2.5cm}p{3.6cm}p{4.2cm}@{}}
    \toprule
    Trục & Vai trò & Phương án đã cân nhắc \\
    \midrule
    Selective & duyệt theo mục tiêu & BFS toàn phần (nổ context) \\
    Traversal & thay 1-hop & PageRank tĩnh (bỏ yếu tố thời gian) \\
    \midrule
    Claim & kiểm từng mệnh đề & LLM judge end-to-end (đắt, không truy được) \\
    Auditor & với context & \\
    \midrule
    Citation & ràng buộc top-3 & siết allowlist (abstain 10\%) $\cdot$ \\
    Gate & kiểm bằng code & chặn cứng (tệ hơn — xem slide hạn chế) \\
    \bottomrule
  \end{tabular}
  \vspace{0.5em}
  \TODO{trục thứ 4 nếu có — hoặc nói rõ chỉ 3}
  \tm{4:30–6:00}
\end{frame}

%% ------------------------------------------------------- 6:00–12:00
\begin{frame}{LIVE DEMO — 6 phút}
  \begin{enumerate}
    \item {\bf 1 task end-to-end} \tm{2'}
      \begin{itemize}
        \item nhập câu hỏi thật $\rightarrow$ dossier hoàn chỉnh
      \end{itemize}
    \item {\bf Mở trace, giải thích} \tm{2'}
      \begin{itemize}
        \item tool call nào, memory đọc gì, gate chặn gì, vì sao viết lại
      \end{itemize}
    \item {\bf 1 ca tấn công có kiểm soát} \tm{1,5'}
      \begin{itemize}
        \item indirect injection trong chunk $\rightarrow$ bị chặn, chỉ log
      \end{itemize}
    \item {\bf 1 hành động được bảo vệ} \tm{0,5'}
      \begin{itemize}
        \item nộp hồ sơ $\rightarrow$ \texttt{requires\_human\_token} $\rightarrow$ audit
      \end{itemize}
  \end{enumerate}
  \vspace{0.3em}
  \warn{Chuẩn bị sẵn: terminal + repo + video dự phòng đúng tag \texttt{final}.
  Nếu API/mạng lỗi: chuyển video, nói rõ \emph{“đây là bản ghi, không phải chạy trực tiếp”}.}
  \tm{6:00–12:00}
\end{frame}

%% ------------------------------------------------------ 12:00–13:30
\begin{frame}{Đánh giá: 100 câu $\times$ 3 run $\times$ 2 cấu hình}
  \begin{columns}[T]
    \begin{column}{0.48\textwidth}
      \textbf{Thiết kế}
      \begin{itemize}
        \item 100 câu, ground truth chuyên viên
        \item C1 = full · C2 = bỏ Selective Traversal
        \item 3 run/cấu hình = \textbf{600 task-run}
        \item cùng model, cùng token budget
      \end{itemize}
      \TODO{chốt baseline ngoài chạy được}
    \end{column}
    \begin{column}{0.48\textwidth}
      \textbf{Tiêu chí (định nghĩa bằng code)}
      \begin{itemize}
        \item presence · validity · trúng căn cứ
        \item cả ba cùng lúc = success
        \item cost/task, latency/task — cùng mẫu số
        \item contamination: bộ công khai, đã dùng trong [1]
      \end{itemize}
    \end{column}
  \end{columns}
  \tm{12:00–13:30}
\end{frame}

%% ------------------------------------------------------ 13:30–15:00
\begin{frame}{Kết quả: 96/100 so với 91/100 của paper}
  \begin{tabular}{@{}lrrr@{}}
    \toprule
    & G4 (C1) & G4 (C2) & Paper [1] \\
    \midrule
    Abstain & 0,0\% & \TODO{} & — \\
    Citation presence & 100,0\% & \TODO{} & — \\
    Citation validity & 99,0\% & \TODO{} & — \\
    Trúng căn cứ & 97,0\% & \TODO{} & — \\
    \midrule
    \textbf{Cả ba cùng lúc} & \good{\textbf{96,0\%}} & \TODO{} & \textbf{91,0\%} \\
    \bottomrule
  \end{tabular}
  \vspace{0.6em}
  \begin{block}{Nói đúng mức}
    Proxy 3 điều kiện, \textbf{không} có semantic equivalence $\Rightarrow$ 96\% là
    \warn{cận trên}. Paper 91\% có judge. Không được nói “vượt về Answer Correctness”.
  \end{block}
  \TODO{thay bằng mean $\pm$ SD qua 3 run + kiểm tra có ngoài 1 SD}
  \tm{13:30–15:00}
\end{frame}

%% ------------------------------------------------------ 15:00–16:00
\begin{frame}{Ablation và chi phí}
  \begin{columns}[T]
    \begin{column}{0.5\textwidth}
      \textbf{Ablation C1 vs C2}\\[0.3em]
      \TODO{bảng: cơ chế bị tách, chênh lệch, có ngoài nhiễu không}
      \vspace{0.5em}
      \begin{block}{Điểm đã sửa so với D2}
        Precision@2 \warn{thấp hơn} paper (34,47\% vs 37--39\%) và khác $k$
        $\rightarrow$ \textbf{chưa so sánh được}. D2 ghi “nhất quán” là \warn{sai}.
      \end{block}
    \end{column}
    \begin{column}{0.46\textwidth}
      \textbf{Cost \& latency}\\[0.3em]
      \TODO{cost/task (USD), latency/task (s), \$/100 câu, ngày chạy, model id}
    \end{column}
  \end{columns}
  \tm{15:00–16:00}
\end{frame}

%% ------------------------------------------------------ 16:00–17:00
\begin{frame}{Bảo mật: 20 ca injection, trước và sau khi vá}
  \small
  \begin{tabular}{@{}llll@{}}
    \toprule
    Nhóm & Số ca & Trước khi vá (1 lần/ca) & Sau khi vá \\
    \midrule
    Direct (A01--A07)  & 7 & 5 chặn · 1 phần (A07) · 1 KQCG (A04) & \good{\TODO{/7}} \\
    Indirect (A08--A19) & 12 & 1 chặn (A18) · 1 lọt (A16) · 10 STAGED & \good{\TODO{/12}} \\
    Direct-API (A20)    & 1 & \warn{1 lọt (A20: token `OK')} & \good{\TODO{/1}} \\
    \bottomrule
  \end{tabular}
  \vspace{0.5em}
  \begin{block}{Lỗ hổng đã tìm thấy trong code của chính nhóm (F1--F5 + A20)}
    \texttt{InputSanitizer} có nhưng \warn{không được gọi} · \texttt{evaluate\_query} chỉ
    chặn đầu vào · \texttt{\_CITE} mù khi xuống dòng (A16) · \texttt{approver\_token}
    chỉ kiểm khác rỗng (A20)
  \end{block}
  \TODO{1 ảnh/log hành động nộp hồ sơ A07: token xác thực $\rightarrow$ audit}
  \tm{16:00–17:00}
\end{frame}

%% ------------------------------------------------------ 17:00–18:30
\begin{frame}{Thay đổi từ D2 và hạn chế}
  \begin{columns}[T]
    \begin{column}{0.5\textwidth}
      \textbf{Đã thay đổi}
      \begin{itemize}
        \item thêm Citation Gate top-3
        \item vá cổng nội dung không tin cậy
        \item sửa kết luận Precision
        \item 1 thí nghiệm thất bại đã ghi lại
      \end{itemize}
    \end{column}
    \begin{column}{0.46\textwidth}
      \textbf{Hạn chế}
      \begin{itemize}
        \item thiếu judge semantic equivalence
        \item 100 câu: mẫu nhỏ, công khai
        \item chưa phủ ảnh scan PDF
        \item \TODO{residual risk: 3 điểm}
      \end{itemize}
    \end{column}
  \end{columns}
  \tm{17:00–18:30}
\end{frame}

%% ------------------------------------------------------ 18:30–20:00
\begin{frame}{Kết luận}
  \begin{itemize}
    \item Giả thuyết: \TODO{bác bỏ / không bác bỏ — kèm số và độ biến thiên}
    \item Hệ thống chạy thật, end-to-end, có trace và audit
    \item Truy xuất vượt paper ở Recall/F2/nhiễu; \warn{thua Precision@2}
    \item Mở: Answer Correctness thật cần LLM-judge; mở rộng ngoài ngành ngân hàng
  \end{itemize}
  \vspace{1em}
  \begin{center}
    \Large Cảm ơn — sẵn sàng hỏi đáp\\[0.3em]
    \normalsize\color{gray}10 phút Q\&A: mỗi thành viên trả lời phần mình phụ trách
  \end{center}
  \tm{18:30–20:00}
\end{frame}

%% =====================================================================
%%  BACKUP — KHÔNG tính vào 20 phút. Dùng khi bị hỏi sâu.
%% =====================================================================
\appendix

\begin{frame}{Backup: kiến trúc tham chiếu đầy đủ}
  \TODO{Hình 1 bản lớn, có tên từng agent + file code tương ứng}
\end{frame}

\begin{frame}{Backup: bảng 20 ca injection chi tiết}
  \TODO{Bảng đầy đủ A01--A20: ID · loại · điểm vào · phòng vệ · kết quả · residual risk}
\end{frame}

\begin{frame}{Backup: trace đầy đủ 1 task}
  \TODO{log thật, đánh số mũi tên}
\end{frame}

\begin{frame}{Backup: bảng đối chiếu yêu cầu R1–R4 + RT1–RT4}
  \TODO{từ Appendix D của báo cáo}
\end{frame}

\begin{frame}{Backup: cost chi tiết theo cấu hình}
  \TODO{token in/out, đơn giá, tổng USD, ngày chạy}
\end{frame}

\begin{frame}{Backup: Tài liệu tham khảo}
  \scriptsize
  \begin{enumerate}
    \item[\textbf{[1]}] K. N. Phan, X.-B. Le, T. T. Quan. \emph{SBV-LawGraph: A Hybrid RAG Approach Integrating Knowledge Graph for the State Bank of Vietnam Legal Documents}. ACIIDS, Springer, 2026.
    \item[\textbf{[2]}] S. Yao et al. \emph{ReAct: Synergizing Reasoning and Acting in Language Models}. ICLR, 2023.
    \item[\textbf{[3]}] A. Asai et al. \emph{Self-RAG: Learning to Retrieve, Generate, and Critique through Self-Reflection}. ICLR, 2024.
    \item[\textbf{[4]}] D. Edge et al. \emph{From Local to Global: A Graph RAG Approach to Query-Focused Summarization}. arXiv:2404.16130, 2024.
    \item[\textbf{[5]}] J. Wei et al. \emph{Long-form Factuality in Large Language Models (SAFE)}. arXiv:2403.18802, 2024.
    \item[\textbf{[6]}] K. Greshake et al. \emph{Compromising Real-World LLM-Integrated Applications with Indirect Prompt Injection}. ACM AISec, 2023.
    \item[\textbf{[7]}] E. Debenedetti et al. \emph{AgentDojo: Evaluating Prompt Injection Attacks and Defenses for LLM Agents}. NeurIPS, 2024.
    \item[\textbf{[8]}] E. Wallace et al. \emph{The Instruction Hierarchy: Training LLMs to Prioritize Privileged Instructions}. arXiv:2404.13208, 2024.
  \end{enumerate}
\end{frame}

\end{document}
